\HLASourcePage{40}{39}
\thispagestyle{hla-body}
\begin{hlafixedpage}
  \HLAFixedBlock{90}{69.79}{432}{\rmfamily}{9.96}{12}{%
    Once a federation execution exists, federates may join and resign from it in any sequence that is meaningful\linebreak
    to the federation user. Figure 2 presents a generalized view of the basic relationship between a federate and\linebreak
    the RTI during the federate participation in a federation execution. The broad arrows in Figure 2 represent\linebreak
    the general invocation of HLA service groups and are not intended to demonstrate strict ordering\linebreak
    requirements on the use of the services. An RTI shall allow an application to participate in a federation as\linebreak
    multiple joined federates, and an RTI shall allow a single application to participate in multiple federations as\linebreak
    a joined federate or federates.%
  }

  \HLAFigureImageAt{138.66}{193.38}{334.68}{251.04}{assets/extracted/source-p040-xref13869.png}
  \HLAFigureCaptionAt{2}{456.17}{Overall view of federate-to-RTI relationship}

  \HLAFixedBlock{90}{508.14}{432}{\rmfamily}{9.96}{12}{%
    The state diagram in Figure 3 is the first of a series of hierarchical state diagrams that formally describe the\linebreak
    exchanges between federates and an RTI. These state diagrams are formal, accurate descriptions of federate\linebreak
    state information depicted in the highly structured, compact, and expressive state chart notation pioneered\linebreak
    by David Harel [B4].\HLASupMark{14} This specification contains two independent sets of state charts. The first set,\linebreak
    beginning in Figure 3 and including all other state charts except Figure 6 through Figure 9, describes the\linebreak
    state of a given federate, from the perspective of that federate, in varying levels of detail. The second set,\linebreak
    Figure 6 through Figure 9, describes the state of a federation execution (including the RTI and all joined\linebreak
    federates) strictly with respect to the use of synchronization points. Finally, Figure 1 is unrelated to either of\linebreak
    these sets of state charts. It is a stand-alone overview of the basic states of a federation execution.%
  }

  \HLAFixedBlock{90}{628.14}{432}{\rmfamily}{9.96}{12}{%
    The next few subclauses describe the first two of these state charts in detail as a way of introducing some of\linebreak
    Harel’s notation and providing an understanding of how the complete collection of state charts in this\linebreak
    document are hierarchically interrelated.%
  }

  \draw[line width=0.48pt] (90.78,702.00)--(127.26,702.00);
  \HLAFixedBlock{90}{704.10}{6}{\rmfamily\raggedright}{6.36}{7}{14}
  \HLAFixedBlock{96.36}{707.25}{425.64}{\rmfamily\raggedright}{7.98}{9.4}{%
    Numbers in brackets correspond to those in the bibliography in Annex I.%
  }
\end{hlafixedpage}
\clearpage
